% 必须使用 XeLaTeX 编译
\documentclass[12pt, UTF8]{ctexart}

% 引入自定义样式包
\usepackage{ructhesis}

% 参考文献设置（保留在主文档便于管理源文件与样式更改）
\usepackage[backend=biber, style=gb7714-2015, sorting=gb7714-2015]{biblatex}
\addbibresource{references.bib}

\begin{document}

% ==========================================
% 封面页
% ==========================================
\begin{titlepage}
\thispagestyle{fancy}
\cfoot{}
  \begin{flushright}
    \heiti \textbf{论文编码：RUC-BK-050101-2018000000}
  \end{flushright}

  \vspace*{2em}
  \begin{center}
      \fontsize{28pt}{\baselineskip}\textbf{\heiti{中国人民大学本科毕业论文（设计）}}
      \vspace{24mm}
      
      % 论文主标题
      \fontsize{28pt}{\baselineskip}\textbf{\heiti 植树造林的跨区域水文溢出与农业经济效应}
      \vspace{3mm}
      
      % 论文副标题
      \begin{spacing}{1.2}
        \LARGE\selectfont{\textbf{\heiti ——来自中国大规模造林工程的证据}}
      \end{spacing}
  \end{center}

  \vspace{40mm}

  \flushleft
  \begin{spacing}{1.3}
    \hspace{27mm}\heiti\LARGE\selectfont{\textbf{作\hspace{14mm}者：}\dunderline[-10pt]{1pt}{\makebox[78mm][c]{庄妍}}}
    
    \hspace{27mm}\heiti\LARGE\selectfont{\textbf{学\hspace{14mm}院：}\dunderline[-10pt]{1pt}{\makebox[78mm][c]{生态环境学院}}}
    
    \hspace{27mm}\heiti\LARGE\selectfont{\textbf{专\hspace{14mm}业：}\dunderline[-10pt]{1pt}{\makebox[78mm][c]{资源与环境经济学}}}
    
    \hspace{27mm}\heiti\LARGE\selectfont{\textbf{年\hspace{14mm}级：}\dunderline[-10pt]{1pt}{\makebox[78mm][c]{2022级}}}
    
    \hspace{27mm}\heiti\LARGE\selectfont{\textbf{学\hspace{14mm}号：}\dunderline[-10pt]{1pt}{\makebox[78mm][c]{}}}
    
    \hspace{27mm}\heiti\LARGE\selectfont{\textbf{指导教师：}\dunderline[-10pt]{1pt}{\makebox[78mm][c]{}}}
    
    \hspace{27mm}\heiti\LARGE\selectfont{\textbf{论文成绩：}\dunderline[-10pt]{1pt}{\makebox[78mm][c]{}}}
    
    \hspace{27mm}\heiti\LARGE\selectfont{\textbf{完成日期：}\dunderline[-10pt]{1pt}{\makebox[78mm][c]{}}}
  \end{spacing}
\end{titlepage}

% ==========================================
% 摘要页
% ==========================================
\newpage
\begin{abstract}
\begin{onehalfspace}
此处撰写中文摘要。重点阐述研究背景、因果推断方法、核心政策评估结果以及空间溢出效应的结论。
\\[12pt]
\textbf{\heiti 关键词：}政策评估 \quad 空间计量 \quad 正外部性
\end{onehalfspace}
\end{abstract}

\setcounter{page}{1}
\pagenumbering{Roman}
\cfoot{\footnotesize \thepage}
\newpage

\begin{enabstract}
\begin{doublespace}
This is the abstract. It summarizes the background, methodology (e.g., causal inference, spatial data analysis), and main results concerning hydrological externalities.
\\[12pt]
\textbf{Keywords：}Policy Evaluation \quad Spatial Econometrics \quad Positive Externalities
\end{doublespace}
\end{enabstract}
\cfoot{\thepage}
\newpage
\cfoot{}

% ==========================================
% 目录页
% ==========================================
\renewcommand\contentsname{内容目录}
\renewcommand\listfigurename{插\ 图}
\renewcommand\listtablename{表\ 格}

\tableofcontents
\listoffigures
\listoftables
\newpage

% ==========================================
% 正文部分
% ==========================================
\setcounter{page}{1}
\pagenumbering{arabic}
\pagestyle{fancy}
\cfoot{\zihao{5} 第 \thepage 页}

\begin{spacing}{1.25}

\section{引言}
在此处输入正文。针对空间数据分析和计量模型的设定\footnote{空间权重矩阵的设定逻辑说明。}，可结合具体的模型展开。

\subsection{研究背景与意义}
具体的内容论述。

\section{文献综述}
文献评述部分。

\end{spacing}

% ==========================================
% 参考文献与附录
% ==========================================
\newpage
\addcontentsline{toc}{section}{\heiti 参考文献}
\printbibliography[heading=bibliography,title=参考文献]

\newpage
\addcontentsline{toc}{section}{\heiti 附录}
\section*{附录}
\appendix
补充的计量回归表格、代码说明或数据来源描述。

\newpage
\addcontentsline{toc}{section}{\heiti 致谢}
\section*{致谢}
致谢部分。

\end{document}